\chapter{一致收敛的保持定理}
\label{chap:uniform-convergence-preservation}

函数列取极限以后，连续性是否保留，积分和导数能否逐项进行，取决于极限
过程能否与邻近点比较、有限求和或差商相协调。一致收敛提供了对所有自变量
共同有效的误差估计，足以处理连续性和 Riemann 积分；求导涉及除以趋于零
的增量，还必须控制导数列。本章逐一证明这些结论，并用反例说明条件缺失时
误差集中在何处。

\section{连续函数的一致极限与局部一致收敛}

\begin{theorem}[一致极限保持连续性]
\label{thm:uniform-limit-continuous}
设 \(E\subseteq\R\)，每个 \(f_n:E\to\R\) 都在 \(x_0\in E\) 连续。
若 \(f_n\rightrightarrows f\) 于 \(E\)，则 \(f\) 在 \(x_0\) 连续。
特别地，连续函数列的一致极限仍连续。
\end{theorem}

\begin{proof}
给定 \(\varepsilon>0\)，由一致收敛取 \(N\)，使
\[
 |f_N(x)-f(x)|<\frac{\varepsilon}{3}
 \qquad(x\in E).
\]
由 \(f_N\) 在 \(x_0\) 连续，存在 \(\delta>0\)，使
\(x\in E\)、\(|x-x_0|<\delta\) 时
\[
 |f_N(x)-f_N(x_0)|<\frac{\varepsilon}{3}.
\]
于是
\[
\begin{aligned}
 |f(x)-f(x_0)|
 &\le |f(x)-f_N(x)|+|f_N(x)-f_N(x_0)|\\
 &\quad+|f_N(x_0)-f(x_0)|<\varepsilon.
\end{aligned}
\]
\end{proof}

证明先固定一个充分靠后的 \(f_N\)，再借助它的连续性比较 \(x\) 与
\(x_0\)。逐点收敛做不到第一步，因为在 \(x\) 与 \(x_0\) 处可能需要
不同的指标。

在非紧区间上，全局一致收敛有时强于局部问题实际所需的条件。

\begin{definition}[局部一致收敛]
设 \(I\subseteq\R\) 为区间。若对每个紧子区间 \(K\subseteq I\)，都有
\[
 f_n\rightrightarrows f\quad\text{于 }K,
\]
则称 \(f_n\) 在 \(I\) 上\term{局部一致收敛}于 \(f\)。
\end{definition}

若 \(I\) 本身是紧区间，局部一致收敛与一致收敛等价。若 \(I\) 非紧，
局部一致收敛只控制每个固定紧子区间，所需指标可以随 \(K\) 改变。

\begin{corollary}
连续函数列的局部一致极限仍连续。
\end{corollary}

\begin{proof}
给定 \(x_0\in I\)，选取含 \(x_0\) 的一个紧相对子区间 \(K\)，使
\(K\) 还含有 \(x_0\) 在 \(I\) 中的某个相对邻域。在 \(K\) 上使用
\theoremref{thm:uniform-limit-continuous}，便得 \(f\) 在 \(x_0\)
连续。
\end{proof}

例如 \(\sum_{k=0}^n x^k\) 在 \((-1,1)\) 上局部一致收敛，因为每个
\([-r,r]\)、\(r<1\) 上的几何尾项都有共同估计；它不在整个
\((-1,1)\) 上一致收敛。

\section{极限与 Riemann 积分交换}

闭区间上的 Riemann 积分满足
\[
 \left|\int_a^b h(x)\,dx\right|
 \le(b-a)\norm h_\infty.
\]
这个估计正好把上确界误差送入积分误差。在应用它之前，还需证明一致极限
本身可积。

\begin{theorem}[一致极限的 Riemann 可积性]
\label{thm:uniform-limit-riemann-integrable}
设每个 \(f_n\) 都在 \([a,b]\) 上 Riemann 可积，并且
\(f_n\rightrightarrows f\)。则 \(f\) Riemann 可积，且
\[
 \int_a^b f(x)\,dx
 =\lim_{n\to\infty}\int_a^b f_n(x)\,dx.
\]
\end{theorem}

\begin{proof}
由一致收敛和某个 \(f_N\) 的有界性，\(f\) 有界。先证可积性。给定
\(\varepsilon>0\)，取 \(N\) 使
\[
 \norm{f-f_N}_\infty<\frac{\varepsilon}{4(b-a)}
\]
（\(a=b\) 时结论直接成立）。因 \(f_N\) 可积，存在分割 \(P\) 使
\[
 U(f_N,P)-L(f_N,P)<\frac{\varepsilon}{2}.
\]
在 \(P\) 的任一子区间 \(I\) 上，
\[
 \sup_I f-\inf_I f
 \le \sup_I f_N-\inf_I f_N+2\norm{f-f_N}_\infty.
\]
乘以各子区间长度并求和，得到
\[
 U(f,P)-L(f,P)
 <\frac{\varepsilon}{2}
   +2(b-a)\frac{\varepsilon}{4(b-a)}
 =\varepsilon.
\]
由 Darboux 判据，\(f\) 可积。

此后 \(f_n-f\) 也可积，并有
\[
 \left|\int_a^b f_n-\int_a^b f\right|
 \le(b-a)\norm{f_n-f}_\infty\to0.
\]
\end{proof}

\begin{corollary}[积分上限函数的一致收敛]
在上述条件下，令
\[
 F_n(x)=\int_a^x f_n(t)\,dt,\qquad
 F(x)=\int_a^x f(t)\,dt.
\]
则 \(F_n\rightrightarrows F\) 于 \([a,b]\)，并且
\[
 \norm{F_n-F}_\infty
 \le(b-a)\norm{f_n-f}_\infty.
\]
\end{corollary}

\begin{proof}
对每个 \(x\in[a,b]\)，
\[
 |F_n(x)-F(x)|
 \le |x-a|\norm{f_n-f}_\infty
 \le(b-a)\norm{f_n-f}_\infty.
\]
对 \(x\) 取上确界即可。
\end{proof}

若只有逐点收敛，上确界误差可能集中在越来越窄或不断移动的区间上，积分
是否收敛便取决于峰高与宽度的共同作用。第五节将给出积分保持不变的尖峰。

\section{极限与微分交换}

一致收敛本身不能控制差商，因为
\[
 \frac{f_n(x+h)-f_n(x)}{h}
\]
把函数误差除以 \(h\)。当 \(h\to0\) 时，再小的固定一致误差也可能被
放大。可行的条件是直接控制导数列。

\begin{theorem}[逐项求导定理]
\label{thm:uniform-derivative-limit}
设 \(f_n:[a,b]\to\R\) 连续并在 \((a,b)\) 上可导。假设：
\begin{enumerate}[label=\textup{(\roman*)}]
 \item 存在 \(x_0\in[a,b]\)，使数列 \((f_n(x_0))\) 收敛；
 \item \(f_n'\) 在 \((a,b)\) 上一致收敛于函数 \(g\)。
\end{enumerate}
则 \((f_n)\) 在 \([a,b]\) 上一致收敛于某个连续函数 \(f\)，并且
\[
 f'(x)=g(x)\qquad(a<x<b).
\]
\end{theorem}

\begin{proof}
先证 \((f_n)\) 一致 Cauchy。对 \(m,n\) 和 \(x\in[a,b]\)，将
Lagrange 中值定理用于 \(f_m-f_n\) 在 \(x\) 与 \(x_0\) 之间的区间，
得到
\[
 |(f_m-f_n)(x)-(f_m-f_n)(x_0)|
 \le(b-a)\sup_{a<t<b}|f_m'(t)-f_n'(t)|.
\]
右端因导数列一致收敛而在 \(m,n\to\infty\) 时趋于零，第一项
\((f_m-f_n)(x_0)\) 也由假设趋于零。因此 \((f_n)\) 满足一致 Cauchy
条件，存在一致极限 \(f\)。由上一节定理，\(f\) 连续。

固定 \(x\in(a,b)\) 和 \(\varepsilon>0\)。导数列一致 Cauchy，所以可取
\(N\)，使对所有 \(n\ge N\) 及 \(t\in(a,b)\)，
\[
 |f_n'(t)-f_N'(t)|<\frac{\varepsilon}{3}.
\]
令 \(n\to\infty\)，先对 \(f_n-f_N\) 使用中值定理，可得当
\(x+h\in(a,b)\) 时
\[
 |(f-f_N)(x+h)-(f-f_N)(x)|
 \le\frac{\varepsilon}{3}|h|.
\]
另一方面，一致极限给出
\[
 |g(x)-f_N'(x)|\le\frac{\varepsilon}{3}.
\]
由 \(f_N\) 在 \(x\) 可导，可取充分小的 \(h\ne0\)，使
\[
 \left|
 \frac{f_N(x+h)-f_N(x)}h-f_N'(x)
 \right|<\frac{\varepsilon}{3}.
\]
三项相加便得
\[
 \left|\frac{f(x+h)-f(x)}h-g(x)\right|<\varepsilon.
\]
故 \(f'(x)=g(x)\)。
\end{proof}

条件 (i) 只固定一个加法常数。若 \(f_n(x)=x+n\)，则 \(f_n'\equiv1\)
一致收敛，但函数列在任何点都不收敛。一个点的收敛与导数的一致控制结合，
通过中值定理传播到整个区间。

若已知每个 \(f_n'\) 连续，还可以从
\[
 f_n(x)=f_n(x_0)+\int_{x_0}^x f_n'(t)\,dt
\]
出发，使用上一节的积分交换定理。这一证明较短，却额外要求导数连续；
上面的中值定理证明只需要导数存在。

\section{Dini 定理}

紧致性有时能把带有单调结构的逐点收敛提升为一致收敛。

\begin{theorem}[Dini 定理]
\label{thm:dini}
设 \(K=[a,b]\)，\(f_n\in C(K)\)，且 \(f_n\) 逐点收敛于连续函数
\(f\)。若对每个 \(x\in K\)，序列 \(f_n(x)\) 关于 \(n\) 单调递增，
或都单调递减，则 \(f_n\rightrightarrows f\) 于 \(K\)。
\end{theorem}

\begin{proof}
只证递减情形。令 \(g_n=f_n-f\)，则 \(g_n\in C(K)\)，
\(g_n\ge0\)，且 \(g_n(x)\downarrow0\)。固定 \(\varepsilon>0\)，定义
\[
 U_n=\{x\in K:g_n(x)<\varepsilon\}.
\]
每个 \(U_n\) 在 \(K\) 中相对开，且
\[
 U_1\subseteq U_2\subseteq\cdots,\qquad
 \bigcup_{n=1}^\infty U_n=K.
\]
由 \(K\) 的紧致性，可从这个开覆盖中取有限子覆盖。由于集合递增，有限
子覆盖中指标最大的一个 \(U_N\) 已覆盖 \(K\)。于是
\[
 0\le f_n(x)-f(x)=g_n(x)<\varepsilon
\]
对所有 \(x\in K\)、\(n\ge N\) 成立，故收敛一致。递增情形对
\(f-f_n\) 使用同一论证。
\end{proof}

定理同时使用了三个条件：定义域紧致，极限连续，收敛关于 \(n\) 单调。
在 \([0,\infty)\) 上，\(x/(n+x)\downarrow0\) 且两端函数均连续，但
收敛不一致；这里缺少紧致性。函数列 \(x^n\) 在 \([0,1]\) 上逐点递减，
但极限在端点不连续，也不一致收敛。若删去单调性，移动峰函数已经给出
连续函数列逐点趋于连续函数却不一致的例子。

\section{逐点收敛下的积分与导数反例}

先构造一个面积不消失的尖峰。在 \([0,1]\) 上令
\[
 s_n(x)=
 \begin{cases}
 n^2x(1-nx),&0\le x\le1/n,\\
 0,&1/n<x\le1.
 \end{cases}
\]
每个 \(s_n\) 连续。固定 \(x>0\) 时，充分大的 \(n\) 有 \(x>1/n\)，
故 \(s_n(x)=0\)；而 \(s_n(0)=0\)。所以 \(s_n\to0\) 逐点成立。但令
\(u=nx\)，得到
\[
 \int_0^1s_n(x)\,dx
 =\int_0^1u(1-u)\,du=\frac16.
\]
峰宽约为 \(1/n\)，峰高约为 \(n\)，两者的乘积保持常数量级。逐点极限
与积分不能交换。

求导的失败即使在原函数一致收敛时也会发生。令
\[
 f_n(x)=\frac{x}{1+n^2x^2}\qquad(-1\le x\le1).
\]
由
\[
 \sup_{x\in\R}\frac{|x|}{1+n^2x^2}=\frac1{2n}
\]
可知 \(f_n\rightrightarrows0\)。然而
\[
 f_n'(x)=\frac{1-n^2x^2}{(1+n^2x^2)^2}
\]
在 \(x=0\) 恒等于 \(1\)，在每个 \(x\ne0\) 处趋于 \(0\)。导数列逐点
收敛于在原点取 \(1\)、其余点取 \(0\) 的函数，它不是极限函数 \(0\)
的导数。这里缺少的正是导数的一致控制。

另一个常见例子是 \(n^{-1}\sin(nx)\rightrightarrows0\)，而导数
\(\cos(nx)\) 不收敛。它说明原函数振幅趋零并不限制振荡频率。对函数列
进行极限运算时，连续性只比较邻近函数值，积分会平均误差，微分却会放大
小尺度变化；三类保持定理的条件因而不同。

\section*{习题}
\addcontentsline{toc}{section}{习题}

\subsection*{A组　连续性与积分}
\begin{enumerate}[label=\textbf{\thechapter.\arabic*.}]
 \exerciseitem{uniform-limit-continuity-epsilon}
 不引用定理名称，直接从定义证明：若 \(f_n\rightrightarrows f\) 于
 \(E\)，且每个 \(f_n\) 在 \(x_0\) 连续，则 \(f\) 在 \(x_0\) 连续。
 \exerciseitem{local-uniform-geometric-series}
 证明几何部分和在 \((-1,1)\) 上局部一致收敛于 \(1/(1-x)\)，但不在
 整个区间上一致收敛。
 \exerciseitem{uniform-integrability-darboux}
 把正文中一致极限可积性的 Darboux 估计逐个子区间写出，并说明为何只需
 选取一个 \(f_N\)。
 \exerciseitem{uniform-integral-varying-upper-limit}
 在正文积分上限函数的条件下，证明
 \[
  \sup_{a\le u\le v\le b}
  \left|\int_u^v(f_n-f)\right|
  \le(b-a)\norm{f_n-f}_\infty.
 \]
 \exerciseitem{pointwise-integral-spike}
 核对正文 \(s_n\) 的连续性、逐点极限、最大值和积分，并说明它为什么
 不可能一致收敛于零。
 \exerciseitem{uniform-integral-with-moving-interval}
 若 \(f_n\rightrightarrows f\) 于 \([a,b]\)，且
 \(a_n,b_n,\alpha,\beta\in[a,b]\)、\(a_n\to\alpha\)、\(b_n\to\beta\)，
 证明
 \[
  \int_{a_n}^{b_n}f_n(x)\,dx\to\int_\alpha^\beta f(x)\,dx.
 \]
\end{enumerate}

\subsection*{B组　逐项求导}
\begin{enumerate}[label=\textbf{\thechapter.\arabic*.},resume]
 \exerciseitem{derivative-theorem-anchor}
 说明逐项求导定理中只需在一个点假设 \(f_n(x_0)\) 收敛，并用
 \(f_n(x)=x+n\) 说明这一条件不能完全删除。
 \exerciseitem{derivative-theorem-uniform-function-convergence}
 在逐项求导定理的假设下证明定量估计
 \[
  \norm{f_m-f_n}_\infty
  \le |f_m(x_0)-f_n(x_0)|
    +(b-a)\norm{f_m'-f_n'}_\infty.
 \]
 \exerciseitem{uniform-functions-pointwise-derivatives}
 完整分析正文 \(f_n(x)=x/(1+n^2x^2)\) 的一致收敛和导数逐点极限。
 \exerciseitem{uniform-derivative-no-pointwise-function-given}
 设 \(f_n'\rightrightarrows g\) 于 \((a,b)\)，且 \(f_n(x_0)\) 收敛。
 证明不必预先假设 \(f_n\) 逐点收敛，正文定理已经能推出一致收敛。
 \exerciseitem{termwise-derivative-series-form}
 设 \(u_n\in C[a,b]\) 且在 \((a,b)\) 可导，若
 \(\sum u_n(x_0)\) 收敛且 \(\sum u_n'\) 一致收敛，证明
 \(\sum u_n\) 一致收敛并可逐项求导。
 \exerciseitem{derivative-limit-darboux}
 在正文逐项求导定理的条件下，说明极限函数 \(g\) 必具有 Darboux 性，
 即使没有假设 \(g\) 连续。
\end{enumerate}

\subsection*{C组　Dini 定理与条件辨析}
\begin{enumerate}[label=\textbf{\thechapter.\arabic*.},resume]
 \exerciseitem{dini-closed-set-proof}
 用闭集
 \[
  F_n=\{x\in K:f_n(x)-f(x)\ge\varepsilon\}
 \]
 和嵌套闭集定理重新证明 Dini 定理的递减情形。
 \exerciseitem{dini-compactness-counterexample}
 证明 \(x/(n+x)\) 在 \([0,\infty)\) 上逐点单调递减到零但不一致收敛。
 \exerciseitem{dini-continuous-limit-counterexample}
 用 \(x^n\) 说明 Dini 定理不能删除极限连续性。
 \exerciseitem{dini-monotonicity-counterexample}
 构造连续函数列，在 \([0,1]\) 上逐点趋于连续函数 \(0\)，但不一致收敛，
 并核对它不满足关于 \(n\) 的逐点单调性。
 \exerciseitem{dini-uniform-supremum-proof}
 设 \(f_n\downarrow f\) 满足 Dini 定理的条件。证明连续函数
 \(g_n=f_n-f\) 的最大值 \(M_n\) 单调趋于零，并指出证明
 \(\lim M_n=0\) 时紧致性出现在哪里。
 \optionalexerciseitem{uniform-limit-homeomorphism-inverses}
 设 \(f_n:[a,b]\to[c,d]\) 为严格递增连续双射，\(f_n\rightrightarrows f\)，
 且 \(f\) 也是严格递增双射。证明 \(f_n^{-1}\rightrightarrows f^{-1}\)。
 提示：使用 \(f^{-1}\) 的一致连续性，不要假设存在统一的导数下界。
\end{enumerate}
